\chapter{Giới thiệu}
\label{ch:gioi-thieu}

\section{Đặt vấn đề}
\label{sec:dat-van-de}

Rối loạn trầm cảm (Major Depressive Disorder -- MDD) là một trong những nguyên nhân hàng đầu gây ra tình trạng khuyết tật và suy giảm chất lượng cuộc sống trên quy mô toàn cầu \cite{pratt1991discriminant}. Theo ước tính của Tổ chức Y tế Thế giới (WHO), có hơn 300 triệu người thuộc mọi lứa tuổi đang phải đối mặt với các mức độ trầm cảm khác nhau. Không chỉ ảnh hưởng nghiêm trọng đến sức khỏe tinh thần, năng suất lao động và mối quan hệ xã hội, trầm cảm mức độ nặng còn là căn nguyên trực tiếp dẫn đến hàng trăm nghìn ca tự tử mỗi năm. Đáng lo ngại hơn, phần lớn những người có triệu chứng trầm cảm ở giai đoạn khởi phát hoặc mức độ nhẹ thường không nhận biết được tình trạng của bản thân, hoặc e ngại sự kỳ thị của xã hội (stigma), dẫn đến việc chậm trễ trong việc tiếp cận các dịch vụ chăm sóc và can thiệp y tế kịp thời.

Trong thực hành lâm sàng và các nghiên cứu dịch tễ học hiện nay, công cụ tiêu chuẩn vàng được sử dụng rộng rãi nhất để sàng lọc và đánh giá mức độ triệu chứng trầm cảm là Thang đo Sức khỏe Bệnh nhân 9 mục (Patient Health Questionnaire-9 -- PHQ-9) do Kroenke và các cộng sự phát triển \cite{kroenke2001phq, spitzer1999validation}. PHQ-9 bao gồm 9 câu hỏi tự báo cáo tương ứng với 9 tiêu chuẩn chẩn đoán trầm cảm trong Sổ tay Chẩn đoán và Thống kê Rối loạn Tâm thần (DSM-IV). Mỗi câu hỏi được tính điểm từ 0 đến 3, tạo nên một thang điểm liên tục dao động từ 0 đến 27 điểm. Dựa vào tổng điểm số, y văn chia mức độ nghiêm trọng thành năm bậc lâm sàng rõ rệt: Tối thiểu (Minimal, 0--4 điểm), Nhẹ (Mild, 5--9 điểm), Vừa (Moderate, 10--14 điểm), Nặng vừa (Moderately severe, 15--19 điểm) và Rất nặng (Severe, 20--27 điểm).

Mặc dù PHQ-9 có giá trị giá trị chẩn đoán cao, phương pháp này vẫn tồn tại những hạn chế cố hữu xuất phát từ bản chất tự báo cáo (self-reported questionnaire). Người tham gia khảo sát có thể phóng đại hoặc che giấu triệu chứng do các yếu tố tâm lý chủ quan, rào cản văn hóa, hoặc hiện tượng thiên kiến xã hội (social desirability bias). Do đó, việc tìm kiếm các chỉ dấu sinh học (biomarkers) khách quan có khả năng phản ánh hoặc hỗ trợ dự đoán mức độ triệu chứng trầm cảm đang trở thành một hướng nghiên cứu mũi nhọn trong y học chính xác và tin sinh học.

Sự bùng nổ của các cuộc điều tra sức khỏe quy mô quốc gia, tiêu biểu là Khảo sát Kiểm tra Sức khỏe và Dinh dưỡng Quốc gia Hoa Kỳ (National Health and Nutrition Examination Survey -- NHANES) do Trung tâm Thống kê Y tế Quốc gia (NCHS) thuộc CDC thực hiện định kỳ hai năm một lần \cite{cdc2014nhanes, cdc2016nhanes}, đã cung cấp một nguồn dữ liệu dạng bảng (tabular health data) khổng lồ. NHANES kết hợp phỏng vấn dịch tễ, thăm khám thể lực chi tiết và các xét nghiệm sinh hóa máu chuẩn hóa trong phòng thí nghiệm lưu động (MEC). Các nghiên cứu y sinh học gần đây đã chỉ ra mối liên hệ phức tạp giữa trầm cảm và các rối loạn chuyển hóa toàn thân, bao gồm phản ứng viêm mãn tính thể nhẹ, rối loạn lipid máu, sự dao động của huyết áp và chỉ số khối cơ thể (BMI).

Học máy (Machine Learning -- ML) và Trí tuệ Nhân tạo (AI) đã chứng minh tiềm năng vượt trội trong việc khai phá các mối quan hệ phi tuyến tính phức tạp giữa các chỉ số sinh lý đa chiều để dự đoán nguy cơ bệnh tật. Tuy nhiên, khi áp dụng Machine Learning vào bài toán dự đoán mức độ trầm cảm trên dữ liệu NHANES, cộng đồng nghiên cứu đang phải đối mặt với ba rào cản kỹ thuật nghiêm trọng:
\begin{enumerate}
    \item \textbf{Mất cân bằng lớp cực kỳ nghiêm trọng (Severe Class Imbalance):} Trong quần thể dân cư nói chung, phần lớn người tham gia khảo sát không có hoặc chỉ có triệu chứng tối thiểu ($C_0$, chiếm 74--76\%), trong khi số lượng bệnh nhân có triệu chứng nặng ($C_4$) chỉ chiếm một tỷ lệ rất nhỏ (khoảng 1.0--1.3\%). Sự chênh lệch này khiến các thuật toán học máy thông thường có xu hướng tối ưu hóa bằng cách dự đoán toàn bộ về lớp đa số để đạt độ chính xác bề nổi (Accuracy) rất cao nhưng hoàn toàn thất bại trong việc nhận diện các ca bệnh thực tế.
    \item \textbf{Sự chồng lấn đặc trưng cao (High Feature Overlap):} Khác với các bệnh lý thực thể có chỉ dấu sinh học đặc hiệu rõ ràng, dữ liệu sinh lý và cận lâm sàng giữa các nhóm triệu chứng kế cận (đặc biệt là giữa người bình thường $C_0$ và người có triệu chứng nhẹ $C_1$) có độ phân tán và chồng lấn rất lớn. Các đường ranh giới phân tách trong không gian đặc trưng trở nên mờ nhạt và khó phân định.
    \item \textbf{Sự suy giảm hiệu năng khi tăng độ chi tiết nhãn (Granularity Degradation):} Đa số các công trình trước đây chỉ dừng lại ở bài toán phân loại nhị phân ($T_2$: Có triệu chứng vs Không có triệu chứng). Khi bài toán được mở rộng sang phân loại đa mức độ ($T_3, T_4, T_5$) nhằm đáp ứng nhu cầu phân tầng rủi ro lâm sàng, hiệu năng của các mô hình học máy bị suy giảm đột ngột và chưa có lời giải thỏa đáng.
\end{enumerate}

Xuất phát từ những thách thức thực tiễn và khoảng trống phương pháp luận nêu trên, đề tài \textbf{“Nghiên cứu khung học máy cho dự đoán mức độ triệu chứng trầm cảm từ dữ liệu sinh lý”} được triển khai nhằm xây dựng và đánh giá toàn diện framework mang tên \textbf{DepSev-ML} (Depression Severity Machine Learning). Cần lưu ý rằng, nghiên cứu này không hướng tới việc thay thế bảng câu hỏi PHQ-9 hay đưa ra chẩn đoán lâm sàng tự động thay thế bác sĩ chuyên khoa tâm thần, mà tập trung vào việc nghiên cứu giới hạn dự đoán của các đặc trưng sinh lý dạng bảng và đề xuất các giải pháp kiến trúc học máy phù hợp nhằm nâng cao độ nhạy cho các nhóm bệnh nhân có nguy cơ cao.

\section{Phát biểu bài toán}
\label{sec:phat-bieu-bai-toan}

Bài toán nghiên cứu của DepSev-ML được phát biểu dưới dạng toán học tổng quát của học có giám sát (Supervised Learning) trên dữ liệu dạng bảng.

Cho tập dữ liệu nghiên cứu gồm $N$ quan sát độc lập:
\begin{equation}
    \mathcal{D} = \{(\mathbf{x}_i, s_i)\}_{i=1}^N
\end{equation}
trong đó:
\begin{itemize}
    \item $\mathbf{x}_i \in \mathcal{X} \subseteq \mathbb{R}^d$ là vectơ đặc trưng $d$ chiều của đối tượng thứ $i$, bao gồm các biến số đo lường sinh lý (nhân trắc học, huyết áp), xét nghiệm cận lâm sàng (hóa sinh, huyết học), thông tin lối sống (hoạt động thể lực, giấc ngủ, hút thuốc, rượu bia) và nhân khẩu học.
    \item $s_i \in \{0, 1, 2, \dots, 27\}$ là tổng điểm PHQ-9 của đối tượng thứ $i$, được tính bằng tổng của 9 câu hỏi khảo sát thành phần:
    \begin{equation}
        s_i = \sum_{j=1}^9 \text{DPQ}_{ij}, \quad \text{với } \text{DPQ}_{ij} \in \{0, 1, 2, 3\}
    \end{equation}
\end{itemize}

Nhãn mục tiêu $y_i$ được sinh ra thông qua một hàm ánh xạ lượng giá $\Phi: \{0, \dots, 27\} \to \mathcal{Y}_K$, trong đó $\mathcal{Y}_K = \{0, 1, \dots, K-1\}$ là không gian đầu ra gồm $K$ mức độ nghiêm trọng rời rạc ($K \in \{2, 3, 4, 5\}$). Nghiên cứu định nghĩa bốn thiết lập độ phân giải mục tiêu (Target Granularity):

\begin{itemize}
    \item \textbf{Thiết lập $T_2$ ($K=2$, Phân loại nhị phân):}
    \begin{equation}
        y^{(T_2)}_i = \Phi_{T_2}(s_i) = \begin{cases} 0 & \text{nếu } s_i \in [0, 4] \quad (\text{Minimal / Triệu chứng tối thiểu}) \\ 1 & \text{nếu } s_i \in [5, 27] \quad (\text{Symptomatic / Có triệu chứng}) \end{cases}
    \end{equation}
    \item \textbf{Thiết lập $T_3$ ($K=3$, Ba mức độ gộp quy ước):}
    \begin{equation}
        y^{(T_3)}_i = \Phi_{T_3}(s_i) = \begin{cases} 0 & \text{nếu } s_i \in [0, 4] \quad (C_0: \text{Minimal / Tối thiểu}) \\ 1 & \text{nếu } s_i \in [5, 9] \quad (C_1: \text{Mild / Nhẹ}) \\ 2 & \text{nếu } s_i \in [10, 27] \quad (C_2^+: \text{Severe conv. / Nặng theo quy ước}) \end{cases}
    \end{equation}
    \item \textbf{Thiết lập $T_4$ ($K=4$, Bốn mức độ):}
    \begin{equation}
        y^{(T_4)}_i = \Phi_{T_4}(s_i) = \begin{cases} 0 & \text{nếu } s_i \in [0, 4] \quad (C_0) \\ 1 & \text{nếu } s_i \in [5, 9] \quad (C_1) \\ 2 & \text{nếu } s_i \in [10, 14] \quad (C_2) \\ 3 & \text{nếu } s_i \in [15, 27] \quad (C_3^+: C_3 \cup C_4) \end{cases}
    \end{equation}
    \item \textbf{Thiết lập $T_5$ ($K=5$, Năm mức độ lâm sàng chuẩn PHQ-9):}
    \begin{equation}
        y^{(T_5)}_i = \Phi_{T_5}(s_i) = \begin{cases} 0 & \text{nếu } s_i \in [0, 4] \quad (C_0: \text{Minimal}) \\ 1 & \text{nếu } s_i \in [5, 9] \quad (C_1: \text{Mild}) \\ 2 & \text{nếu } s_i \in [10, 14] \quad (C_2: \text{Moderate}) \\ 3 & \text{nếu } s_i \in [15, 19] \quad (C_3: \text{Moderately severe}) \\ 4 & \text{nếu } s_i \in [20, 27] \quad (C_4: \text{Severe}) \end{cases}
    \end{equation}
\end{itemize}

Mục tiêu học máy là tìm kiếm một hàm dự đoán $f: \mathcal{X} \to \mathcal{Y}_K$ (hoặc hàm ước lượng phân phối xác suất hậu nghiệm $P(Y = k \mid \mathbf{x})$) sao cho tối đa hóa độ đo hiệu năng trên toàn bộ các lớp. Do dữ liệu bị mất cân bằng lớp nghiêm trọng, hàm mục tiêu đánh giá chính của nghiên cứu không sử dụng Accuracy thông thường mà được xác định thông qua \textbf{Macro-F1} và \textbf{Balanced Accuracy}:
\begin{equation}
    \text{Macro-F1} = \frac{1}{K} \sum_{k=0}^{K-1} \text{F1}_k = \frac{1}{K} \sum_{k=0}^{K-1} \frac{2 \cdot \text{Precision}_k \cdot \text{Recall}_k}{\text{Precision}_k + \text{Recall}_k}
\end{equation}
\begin{equation}
    \text{Balanced Accuracy} = \frac{1}{K} \sum_{k=0}^{K-1} \text{Recall}_k
\end{equation}

\textbf{Ràng buộc kỹ thuật bắt buộc (Zero Data Leakage):} Để bảo đảm tính khách quan và khả năng khái quát hóa y sinh thực chất, ma trận đặc trưng đầu vào $\mathbf{X}$ tuyệt đối không được chứa bất kỳ câu hỏi nào trong số 9 câu hỏi khảo sát DPQ. Toàn bộ các bước chuẩn hóa, điền khuyết (imputation) và mã hóa đặc trưng bắt buộc phải được học độc lập bên trong từng tập huấn luyện của kiểm định chéo (Cross-Validation fold) và chỉ được áp dụng lên tập kiểm tra.

\section{Mục tiêu và câu hỏi nghiên cứu}
\label{sec:muc-tieu-va-cau-hoi}

Mục tiêu tổng quát của đề tài là xây dựng một khung học máy thực nghiệm hoàn chỉnh, có khả năng tái lập và kiểm soát chặt chẽ, nhằm khảo sát và nâng cao hiệu năng dự đoán mức độ triệu chứng trầm cảm từ dữ liệu sinh lý đa chu kỳ NHANES. 

Để đạt được mục tiêu tổng quát, nghiên cứu tập trung giải quyết năm câu hỏi khoa học trọng tâm (Research Questions -- RQ):

\begin{itemize}
    \item \textbf{RQ1 (Đóng góp của các nhóm đặc trưng):} Những nhóm đặc trưng nào (trong số nhân trắc học, huyết áp, hóa sinh máu, huyết học, lối sống và nhân khẩu học) đóng góp thực chất nhất vào khả năng dự đoán mức độ PHQ-9, và sự kết hợp đa tầng đặc trưng có giúp cải thiện hiệu năng so với từng nhóm đơn lẻ không?
    \item \textbf{RQ2 (Quy luật suy giảm theo độ chi tiết nhãn):} Hiệu năng của các mô hình học máy thay đổi như thế nào khi độ chi tiết của nhãn mục tiêu tăng dần từ 2 lên 3, 4 và 5 mức độ? Bước chuyển dịch nào và ranh giới nhãn nào là nguyên nhân chính gây ra sự suy giảm hiệu năng lớn nhất?
    \item \textbf{RQ3 (Tác động của việc mở rộng dữ liệu đa chu kỳ):} Liệu việc mở rộng cỡ mẫu từ một chu kỳ đơn lẻ ($N = 5.054$) lên ba chu kỳ gộp ($N = 14.581$) với không gian đặc trưng hài hòa có thể khắc phục được hiện tượng suy giảm hiệu năng ở các mức độ chi tiết hay không?
    \item \textbf{RQ4 (Hiệu quả của phân loại phân cấp):} Kiến trúc phân loại phân cấp hai giai đoạn (Hierarchical Classification) có cải thiện được hiệu năng phân loại trên không gian 3 lớp ($T_3$) so với phân loại phẳng truyền thống (Flat Classification) hay không, và cơ chế định tuyến cứng (Hard Routing) gây ra những hệ quả nào đối với sự lan truyền sai số (Error Propagation)?
    \item \textbf{RQ5 (Khả năng của Soft, Cost-Sensitive và Ordinal Classification):} Việc chuyển đổi từ định tuyến cứng sang định tuyến mềm (Soft Probability Fusion), tối ưu hóa rủi ro Bayes theo ma trận chi phí (Cost-Sensitive Learning) và mô hình hóa trực tiếp quan hệ thứ bậc (Ordinal Classification) giải quyết được những mâu thuẫn nào giữa Macro-F1, Balanced Accuracy, Class MAE và tỷ lệ bỏ sót ca bệnh nặng?
\end{itemize}

\section{Phạm vi nghiên cứu}
\label{sec:pham-vi-nghien-cuu}

Để định hình rõ ranh giới học thuật của đề tài, các giới hạn nghiên cứu được xác định minh bạch như sau:
\begin{enumerate}
    \item \textbf{Miền dữ liệu:} Nghiên cứu giới hạn trên tập dữ liệu khảo sát cắt ngang dạng bảng của NHANES (Hoa Kỳ) ở các chu kỳ 2011--2012, 2013--2014 và 2015--2016, áp dụng cho đối tượng người trưởng thành từ 20 tuổi trở lên có đầy đủ 9 mục khảo sát PHQ-9 hợp lệ.
    \item \textbf{Bản chất tín hiệu:} Đề tài chỉ sử dụng các đặc trưng sinh lý và cận lâm sàng dạng bảng được trích xuất từ hồ sơ khám sức khỏe định kỳ. Nghiên cứu \textbf{không sử dụng tín hiệu chuỗi thời gian thô} như điện tâm đồ (ECG) liên tục hay điện não đồ (EEG).
    \item \textbf{Tính chất bài toán:} Đề tài tập trung vào bài toán lượng giá mức độ triệu chứng trầm cảm tự báo cáo (symptom severity prediction) trong nghiên cứu dịch tễ, \textbf{không phải là hệ thống chẩn đoán bệnh học lâm sàng tự động}. Kết quả đầu ra là điểm số và nhãn nhóm nguy cơ, không thay thế quy trình thăm khám tâm thần chuyên sâu.
    \item \textbf{Đóng góp thuật toán:} Nghiên cứu không tuyên bố phát minh một thuật toán học máy hoàn toàn mới từ gốc, mà tập trung vào việc nghiên cứu kiến trúc hệ thống, tích hợp các phương pháp phân loại phân cấp, quyết định nhạy chi phí và phân loại thứ bậc để giải quyết bài toán mất cân bằng lớp và chồng lấn đặc trưng trong dữ liệu y tế thực tế.
\end{enumerate}

\section{Đóng góp của nghiên cứu}
\label{sec:dong-gop-nghien-cuu}

Những đóng góp khoa học và kỹ thuật chủ yếu của khóa luận bao gồm:
\begin{enumerate}
    \item \textbf{Thiết lập quy trình thực nghiệm chuẩn tắc chống rò rỉ dữ liệu (Strict Zero Leakage Protocol):} Xây dựng khung làm việc DepSev-ML hoàn toàn khép kín, phân tách minh bạch các tập huấn luyện, kiểm định nội bộ và kiểm tra ngoài, mã hóa toàn bộ đường ống tiền xử lý bên trong từng fold, giải quyết triệt để các nguy cơ rò rỉ thông tin trong tiền xử lý và tinh chỉnh siêu tham số.
    \item \textbf{Hài hòa hóa dữ liệu NHANES đa chu kỳ (Data Harmonization):} Xây dựng thành công bộ đặc trưng chuẩn $F_{4\_\text{common31}}$ gồm 31 biến số sinh lý và nhân khẩu học đồng nhất trên 14.581 đối tượng qua 3 chu kỳ NHANES liên tiếp, làm cơ sở thực nghiệm vững chắc cho các đánh giá tổng quát hóa.
    \item \textbf{Phát hiện quy luật suy giảm hiệu năng theo độ chi tiết nhãn (Target Granularity Analysis):} Lượng hóa chính xác sự sụt giảm của Macro-F1 từ 0.6639 ($T_2$) xuống 0.2630 ($T_5$), đồng thời chứng minh bằng thực nghiệm rằng ranh giới giữa người bình thường ($C_0$) và người có triệu chứng nhẹ ($C_1$) là "điểm nghẽn" suy giảm lớn nhất do sự chồng lấn phân phối sinh lý.
    \item \textbf{Phân tích bản chất và khắc phục sự lan truyền sai số trong phân loại phân cấp:} Chỉ ra hạn chế cố hữu của cơ chế định tuyến cứng (Hard Routing) khi loại bỏ nhầm 120 ca nặng quy ước ($C_2^+$) ở giai đoạn đầu; đề xuất mô hình Soft Probability Fusion kết hợp tối ưu chi phí Bayes ($E_{6a}, E_{6b}$) giúp giảm số ca bỏ sót ca nặng xuống còn 62 ca.
    \item \textbf{Hiện thực và chứng minh ưu thế của Ordinal Logistic Regression:} Xây dựng mô hình Proportional Odds với ràng buộc tham số mũ đảm bảo tính đơn điệu của ngưỡng cắt, đạt mức Balanced Accuracy cao nhất \textbf{0.5151} trên Legacy Cohort và duy trì \textbf{0.5073} trên tập kiểm tra theo dòng thời gian (Temporal Evaluation, $N = 4.888$), nâng độ nhạy nhận diện nhóm nặng lên \textbf{60.37\%}.
\end{enumerate}

\section{Cấu trúc báo cáo}
\label{sec:cau-truc-bao-cao}

Nội dung báo cáo được tổ chức thành 5 chương chính:
\begin{itemize}
    \item \textbf{Chương 1: Giới thiệu.} Trình bày bối cảnh nghiên cứu, phát biểu bài toán toán học, mục tiêu, câu hỏi nghiên cứu, phạm vi, đóng góp và cấu trúc báo cáo.
    \item \textbf{Chương 2: Cơ sở lý thuyết và các nghiên cứu liên quan.} Phân tích thang đo PHQ-9, học máy trên dữ liệu dạng bảng, các thuật toán nền tảng (LR, DT, RF, XGBoost), lý thuyết phân loại phân cấp, quyết định nhạy chi phí, phân loại thứ bậc và tổng hợp các công trình liên quan.
    \item \textbf{Chương 3: Phương pháp nghiên cứu và thiết kế DepSev-ML.} Mô tả chi tiết kiến trúc DepSev-ML, quy trình xử lý dữ liệu Legacy và Multi-cycle, không gian nhãn $T_2$--$T_5$, nhóm đặc trưng $F_1$--$F_4$, quy trình kiểm định chéo, thiết kế phân loại phân cấp và thứ bậc, cùng các thuật toán giải thuật giả.
    \item \textbf{Chương 4: Thực nghiệm và đánh giá.} Trình bày chi tiết môi trường thực nghiệm, kết quả Feature Ablation, Nested CV, Target Granularity, Multi-cycle Generalization, so sánh toàn diện các mô hình $E_0$--$E_7$, phân tích chuyên sâu nguyên nhân hiệu năng còn thấp và thảo luận trả lời các câu hỏi RQ1--RQ5.
    \item \textbf{Chương 5: Kết luận và hướng phát triển.} Tổng kết các kết quả đạt được, đóng góp nghiên cứu, các hạn chế khoa học trung thực và đề xuất hướng phát triển ưu tiên tiếp theo.
\end{itemize}

\vspace{0.5cm}
\noindent \textit{Chuyển tiếp sang Chương 2:} Để thiết lập nền tảng khoa học vững chắc phục vụ việc xây dựng khung phương pháp và diễn giải các kết quả thực nghiệm, Chương 2 tiếp theo sẽ trình bày có hệ thống cơ sở lý thuyết về thang đo triệu chứng PHQ-9, các đặc thù toán học của bài toán học máy mất cân bằng trên dữ liệu dạng bảng, cùng các mô hình phân loại đa lớp, phân loại phân cấp và phân loại thứ bậc.
\cleardoublepage
